\section{Conclusiones}

\begin{frame}[c]{Conclusiones}
    \begin{enumerate}
        \item \textbf{El estimador robusto de snipar incrementa el tamaño de muestra efectivo} respecto al método de \textcite{young2022-fgwas}, mejorando la precisión de las estimaciones de heredabilidad en el diseño familiar.
        \item \textbf{Las heredabilidades estimadas en el diseño familiar tienden a ser menores} que en el poblacional. Tres rasgos biomédicos ---diabetes reportada, colesterol HDL y presión arterial diastólica--- muestran el patrón opuesto.
        \item \textbf{El poder estadístico en MCPS es limitado para detectar correlaciones} entre efectos poblacionales ($\hat{\beta}$) y familiares ($\hat{\delta}$), particularmente en rasgos con baja heredabilidad.
        % Esto implica que las correlaciones observadas deben interpretarse a la luz de simulaciones de poder, y no necesariamente como evidencia de ausencia de concordancia entre diseños.
    \end{enumerate}
\end{frame}

\begin{frame}[c]{Perspectivas y trabajo futuro}
    \begin{itemize}
        \item \textbf{Comprobar diferencias de heredabilidad} con pruebas estadísticas apropiadas
        \item \textbf{Complementar comparaciones} con pruebas de heterogeneidad DGEs vs. efectos poblacionales en SNPs indpendientes (\textit{clumping} de $r^2 < 1\%$), con FDR a 5\% y 10\%.
        \item (Opcional) \textbf{Diagnosticar el estimador (correlaciones)} de \textcite{young2022-fgwas}, ¿Cuál es la razón por la que los rasgos de MCPS parece tener bajo poder?
        \item (Opcional) Proceder a fine-mapping.
    \end{itemize}
\end{frame}